%=====================================================================
\chapter{Mapping an ER Diagram to Tables}\label{ch:mapping}
%=====================================================================
\locator{3}{Part III --- From Diagram to Schema}

\begin{outcomes}
By the end of this chapter you will be able to:
\begin{tight}
  \item \textbf{Apply} the seven mapping rules, in order, to any ER diagram.
  \item \textbf{Map} a 1:N relationship by placing a foreign key, and
        \textbf{say} which side it goes on and why.
  \item \textbf{Resolve} an M:N relationship into a table, and
        \textbf{determine} its primary key.
  \item \textbf{Map} a weak entity, a multi-valued attribute and a recursive
        relationship, each of which has its own rule.
  \item \textbf{Check} a finished schema back against the diagram it came
        from, constraint by constraint.
\end{tight}
\end{outcomes}

\begin{prereq}
\chref{cardinality}: the finished, constrained diagram --- you will map the one
you drew. \chref{constraints} for foreign keys and referential actions, which
are what the rules produce.
\end{prereq}

\begin{keyterms}
mapping $\bullet$ logical design $\bullet$ strong entity $\bullet$ weak entity
$\bullet$ relationship table $\bullet$ junction table $\bullet$ composite
primary key $\bullet$ foreign key placement $\bullet$ resolving a
many-to-many $\bullet$ surrogate key $\bullet$ recursive relationship
$\bullet$ ternary relationship $\bullet$ flattening a composite attribute
\end{keyterms}

\section{The Part of Design That Is Mechanical}

Everything in Part~II was judgement: whether a noun is an entity, whether
participation is really total, which candidate key to choose. Those decisions
were expensive and nobody could make them for you.

This chapter is the opposite. Given a finished diagram, turning it into tables
is a \emph{procedure}. Seven rules, applied in order, and two people who
follow them get the same schema. That is the point of having spent four
chapters on the diagram: the hard thinking is over, and what remains can be
checked.

\begin{conceptbox}[The seven rules, in the order they must be applied]
\begin{tabular}{@{}R{0.7cm}L{5.3cm}L{8.0cm}@{}}
\textbf{1} & Strong entity type & one table; its key becomes the primary key \\
\textbf{2} & Composite attribute & store the \emph{parts}, not the whole \\
\textbf{3} & Multi-valued attribute & \textbf{a new table}: owner's key +
  the value, both in the primary key \\
\textbf{4} & 1:1 relationship & a foreign key on the side with \emph{total}
  participation \\
\textbf{5} & 1:N relationship & a foreign key on the \textbf{many} side \\
\textbf{6} & M:N relationship & \textbf{a new table}: both keys, together the
  primary key, plus the relationship's own attributes \\
\textbf{7} & Weak entity type & a table whose primary key is the owner's key
  \emph{plus} the partial key \\
\end{tabular}

\smallskip
\textbf{Order matters.} Rules 4--7 place foreign keys, and a foreign key can
only point at a primary key that already exists. Do all the entities first.
\end{conceptbox}

\section{Rules 1 to 3: the Entities}

\begin{examplebox}[Rule 1 --- a strong entity becomes a table]
\ent{student}, with key \texttt{roll\_no}, becomes

\medskip
\relation{student}{\relkey{roll\_no}, name, programme, semester}

\medskip
\begin{lstlisting}[language=sql]
CREATE TABLE student (
  roll_no   TEXT    PRIMARY KEY NOT NULL,
  name      TEXT    NOT NULL,
  programme TEXT    NOT NULL,
  semester  INTEGER NOT NULL CHECK (semester BETWEEN 1 AND 8)
) STRICT;
\end{lstlisting}

The \sqlkw{CHECK} is the domain constraint from \chref{constraints}, and it
came from the ER diagram's domain, not from inspiration.
\end{examplebox}

\begin{examplebox}[Rules 2 and 3 --- the two attribute kinds that will not fit]
\textbf{Rule 2, composite.} \texttt{address} = house, street, city. Store the
parts:

\begin{lstlisting}[language=sql]
  addr_house  TEXT,
  addr_street TEXT,
  addr_city   TEXT,
\end{lstlisting}

Three columns, one for each leaf of the ER tree. Never one \texttt{address}
column --- \chref{ermodel} explained what that costs, and the cost is paid at
the moment somebody asks ``how many students live in Bahawalpur''.

\medskip
\textbf{Rule 3, multi-valued.} \texttt{phone} cannot be a column at all. It
becomes its own table:

\medskip
\relation{student\_phone}{\relkey{roll\_no}, \relkey{phone}}

\begin{lstlisting}[language=sql]
CREATE TABLE student_phone (
  roll_no TEXT NOT NULL REFERENCES student(roll_no) ON DELETE CASCADE,
  phone   TEXT NOT NULL,
  PRIMARY KEY (roll_no, phone)
) STRICT;
\end{lstlisting}

\smallskip
\textbf{Both attributes are in the primary key, and that is the rule, not a
convenience.} It says a student may have many numbers and may not have the
same number recorded twice. Make \texttt{roll\_no} alone the key and each
student gets one phone --- which is where you started. Leave the key off
altogether and \chref{relational}'s step 2 happens: the same number goes in
twice and no \sqlkw{DELETE} can remove one of them.

\smallskip
The cascade is right here because a phone number belongs to a student in
exactly the way \chref{constraints} requires: it cannot exist without one.
\end{examplebox}

\section{Rules 4 to 6: the Relationships}

\begin{definitionbox}[Where a foreign key goes, and why there is no choice]
In a \textbf{1:N} relationship the foreign key goes on the \textbf{many} side,
always. The reason is the atomicity of \chref{relational}: a course belongs to
one department, so one value fits in one column. A department has many courses,
and many values do not.

\smallskip
In a \textbf{1:1} relationship either side would work, so there is a choice and
a rule for making it: \textbf{put it on the side with total participation}, so
that it can be declared \sqlkw{NOT NULL}. Put it on the partial side and most
of the column is null.

\smallskip
In an \textbf{M:N} relationship \emph{neither} side works, and that is why the
relationship needs a table of its own.
\end{definitionbox}

\begin{worked}[Rule 6 --- resolving \ent{enrols}, and finding its key]
\ent{student} (0,N) --- \ent{enrols} --- (0,N) \ent{section}, carrying
\texttt{grade}.

\medskip
\textbf{Step 1 --- why no foreign key will do.} Put \texttt{section\_id} on
\ent{student} and a student takes one section. Put \texttt{roll\_no} on
\ent{section} and a section holds one student. Both are wrong, and neither is
fixable by choosing the other side.

\medskip
\textbf{Step 2 --- make a table.} It gets the primary key of each participant:

\medskip
\relation{enrolment}{\relkey{roll\_no}, \relkey{course\_code},
\relkey{semester}, \relkey{section\_no}, grade}

\medskip
Four key attributes, not two --- because \ent{section} is weak
(\chref{cardinality}) and its own key is already three attributes wide. A weak
entity on either end of an M:N relationship makes the junction table wide, and
that is the honest consequence of the design rather than a mistake.

\medskip
\textbf{Step 3 --- the primary key is the combination.} All four together.
This states exactly one thing, and it is the thing the university means:
\emph{a student may appear in a section only once}.

\medskip
\textbf{Step 4 --- the relationship's attribute comes too.} \texttt{grade} is
not in the key. It is a fact \emph{about} the enrolment, nullable because the
semester is not over.

\medskip
\begin{lstlisting}[language=sql]
CREATE TABLE enrolment (
  roll_no     TEXT    NOT NULL REFERENCES student(roll_no),
  course_code TEXT    NOT NULL,
  semester    TEXT    NOT NULL,
  section_no  INTEGER NOT NULL,
  grade       TEXT    CHECK (grade IN ('A','B','C','D','F')),
  PRIMARY KEY (roll_no, course_code, semester, section_no),
  FOREIGN KEY (course_code, semester, section_no)
    REFERENCES section(course_code, semester, section_no)
) STRICT;
\end{lstlisting}

\smallskip
\textbf{Notice the second foreign key is written separately.} A foreign key of
three columns cannot be declared inline against one column; it needs its own
\sqlkw{FOREIGN KEY} clause naming all three, in the same order as the parent's
primary key. This is where a wide weak-entity key is actually felt, and it is
the strongest practical argument for giving \ent{section} a surrogate
\texttt{section\_id} instead --- an argument \chref{evolution} makes with a
measurement rather than with taste.
\end{worked}

\begin{alertbox}[The M:N table is not ``just a link table''; it is where the facts live]
Students meet this table and call it plumbing. It is the most informative
table in the database:

\begin{tight}
  \item \texttt{grade} is there, and nowhere else. Every transcript, every GPA,
        every pass rate is a query over this table.
  \item It is the only place that records \emph{when} a student took something.
  \item It is almost always the largest table --- 4{,}000 students $\times$ 5
        sections $\times$ 8 semesters is 160{,}000 rows against a few thousand
        everywhere else, which is why \chref{indexes} measures this table and
        not the others.
\end{tight}

\smallskip
If an M:N table has no attributes of its own today, it still will: ``when did
they enrol'', ``who approved it'', ``was it a repeat''. Design it as a first
class table, because it is one.
\end{alertbox}

\dsfig{%
\begin{tikzpicture}[
  ent/.style={draw=primary, line width=1pt, fill=primary!8,
              text=primary!60!black, font=\scriptsize\bfseries,
              minimum width=1.8cm, minimum height=0.62cm},
  weak/.style={ent, draw=purplehl, fill=purplehl!8, text=purplehl!65!black,
               double, double distance=1.1pt},
  rel/.style={draw=greenhl, line width=0.9pt, fill=greenhl!8,
              text=greenhl!55!black, font=\scriptsize\bfseries,
              diamond, aspect=2.2, inner sep=0.5pt},
  tb/.style={rounded corners=2pt, draw=#1, line width=0.9pt, fill=#1!6,
             text=#1!60!black, font=\scriptsize, align=left,
             inner sep=3pt, text width=5.2cm},
  ln/.style={draw=neutral!55, line width=0.7pt},
  rr/.style={-{Stealth[length=2.2mm]}, accent!75, line width=1pt},
  rl/.style={font=\scriptsize\bfseries, text=accent},
  hd/.style={font=\scriptsize\bfseries, text=primary}
]
  \node[hd] at (-3.4,2.5) {the diagram};
  \node[hd] at ( 4.4,2.5) {the tables};

  \node[ent]  (stu) at (-5.3, 1.3) {STUDENT};
  \node[rel]  (enr) at (-3.4, 1.3) {ENROLS};
  \node[weak] (sec) at (-1.5, 1.3) {SECTION};
  % Double diamond: HAS is the identifying relationship, and drawing it as an
  % ordinary one here contradicted Figure 6.1 two chapters earlier.
  \node[rel, draw=purplehl, fill=purplehl!8, text=purplehl!65!black,
        double, double distance=1.1pt] (has) at (-1.5,-0.3) {HAS};
  \node[ent]  (crs) at (-1.5,-1.6) {COURSE};
  \draw[ln] (stu) -- (enr);  \draw[ln] (enr) -- (sec);
  \draw[ln, double, double distance=1pt] (sec) -- (has);
  \draw[ln] (has) -- (crs);

  \node[tb=primary]  (t1) at (4.4, 1.75)
       {\textbf{student}(\relkey{roll\_no}, name, \dots)};
  \node[tb=accent]   (t2) at (4.4, 0.75)
       {\textbf{enrolment}(\relkey{roll\_no}, \relkey{course\_code},
        \relkey{semester}, \relkey{section\_no}, grade)};
  \node[tb=purplehl] (t3) at (4.4,-0.55)
       {\textbf{section}(\relkey{course\_code}, \relkey{semester},
        \relkey{section\_no}, room)};
  \node[tb=primary]  (t4) at (4.4,-1.6)
       {\textbf{course}(\relkey{course\_code}, title, \dots)};

  % Each label rides its OWN curve with pos=, rather than floating between
  % two of them: the first draft put "rule 6" in the gap where the enrolment
  % and section arrows cross, and it was impossible to tell which it named.
  \draw[rr] (stu.east) ..controls (-4.2,2.6) and (0.6,2.9)..
       node[rl, above, pos=0.55] {rule 1} (t1.west);
  % These two curves run close together in the middle of the figure, so each
  % label goes where its OWN curve is alone: rule 6 late, on the flat approach
  % to the enrolment table; rule 7 early, just as it leaves SECTION. Placed at
  % mid-curve they printed on each other's arrows.
  \draw[rr] (enr.south) ..controls (-3.4,0.2) and (0.6,0.75)..
       node[rl, above, pos=0.80] {rule 6} (t2.west);
  \draw[rr] (sec.east) ..controls (0.0,1.3) and (0.6,-0.2)..
       node[rl, right, pos=0.18] {rule 7} (t3.west);
  \draw[rr] (crs.east) ..controls (0.0,-1.6) and (0.6,-1.6)..
       node[rl, above, pos=0.5] {rule 1} (t4.west);
\end{tikzpicture}}%
{Four of the seven rules on the part of the diagram that uses them. Rule~1
makes a table per box; rule~6 turns the diamond into a table of its own;
rule~7 gives the weak entity its owner's key. The arrows are the whole of
logical design.}{mapping}

\section{Rule 7, and the Two Awkward Cases}

\begin{examplebox}[Rule 7 --- a weak entity]
\ent{section} is identified by its owner plus its discriminator, so its primary
key is the owner's key plus the partial key, and the owner's key is
\emph{simultaneously} a foreign key:

\begin{lstlisting}[language=sql]
CREATE TABLE section (
  course_code TEXT    NOT NULL
              REFERENCES course(course_code) ON DELETE CASCADE,
  semester    TEXT    NOT NULL,
  section_no  INTEGER NOT NULL CHECK (section_no > 0),
  room        TEXT,
  emp_id      TEXT REFERENCES instructor(emp_id) ON DELETE SET NULL,
  PRIMARY KEY (course_code, semester, section_no)
) STRICT;
\end{lstlisting}

\smallskip
Three things are true of \texttt{course\_code} at once: it is part of the
primary key, it is a foreign key, and it is \sqlkw{NOT NULL}. All three come
straight off the diagram --- the double line in Figure~\ref{fig:ercomplete}
said the participation was total, and \sqlkw{CASCADE} is correct here for the
reason \chref{constraints} gave and nowhere else in this schema.
\end{examplebox}

\begin{conceptbox}[The recursive relationship, and the ternary one]
Two cases the seven rules cover, but which look as though they need an eighth.

\medskip
\textbf{Recursive (\ent{course} \ent{requires} \ent{course}, M:N).} Rule 6
applies unchanged. The new table takes the key of each participant --- and
both participants are \ent{course}, so it takes \texttt{course\_code}
\emph{twice}, under the two role names:

\begin{lstlisting}[language=sql]
CREATE TABLE prerequisite (
  course_code TEXT NOT NULL REFERENCES course(course_code),
  requires    TEXT NOT NULL REFERENCES course(course_code),
  PRIMARY KEY (course_code, requires),
  CHECK (course_code <> requires)          -- nothing requires itself
) STRICT;
\end{lstlisting}

The role names from \chref{ermodel} become the column names. Without them you
would have two columns called \texttt{course\_code} and no legal way to name
either.

\medskip
\textbf{Ternary.} A relationship among three entity types becomes one table
with three foreign keys, and its primary key is normally all three. Do
\emph{not} map it as three binary relationships: ``student, section,
instructor'' as three pairs cannot record which instructor taught which
student in which section, only that each pair occurred.
\end{conceptbox}

\begin{pitfall}
\begin{tight}
  \item \textbf{Putting the foreign key on the ``one'' side of a 1:N.} One
        column cannot hold many values. This is the single most common error.
  \item \textbf{Mapping M:N with a foreign key.} It cannot be done; make the
        table.
  \item \textbf{Giving the M:N table a surrogate key and nothing else.} An
        \texttt{id} column makes rows addressable; without the composite
        \sqlkw{UNIQUE} the same student enrols twice.
  \item \textbf{Making the multi-valued table's key just the owner.} Then the
        owner has one value again, which is what you were escaping.
  \item \textbf{Storing a composite attribute whole.} Three columns, one per
        leaf.
  \item \textbf{Declaring a multi-column foreign key inline.} It needs its own
        \sqlkw{FOREIGN KEY} clause, with the columns in the parent's key order.
  \item \textbf{Mapping a ternary relationship as three binary ones.} It loses
        exactly the fact the relationship existed to record.
\end{tight}
\end{pitfall}

\begin{lab}[Turn your diagram into the database you will use all term]
This lab produces \texttt{code/schema/v1\_create.sql}, and every later chapter
runs against what you build here.

\begin{lstlisting}[language=bash]
# --- 1. Start the file ----------------------------------------------
# Open code/schema/v1_create.sql in an editor. First line, always:
#   PRAGMA foreign_keys = ON;
\end{lstlisting}
\begin{lstlisting}[language=sql]
-- --- 2. Rule 1: one table per strong entity ---------------------
--     department, course, instructor, student.
--     Primary key first, NOT NULL, STRICT on every table.

-- --- 3. Rules 2 and 3: the awkward attributes -------------------
--     Flatten the composite one into its parts.
--     Give the multi-valued one its own table, with BOTH
--     attributes in the primary key.

-- --- 4. Rule 5: the 1:N foreign keys ----------------------------
--     course -> department.  Which side, and NOT NULL or not?
--     Your Chapter 6 diagram already decided both.

-- --- 5. Rule 7: the weak entity ---------------------------------
--     section. Owner's key plus discriminator, and the owner's
--     key is a foreign key at the same time.

-- --- 6. Rule 6: the M:N tables ----------------------------------
--     enrolment, and prerequisite (the recursive one).
--     Composite primary keys. Relationship attributes included.

-- --- 7. A referential action for every foreign key --------------
--     There are at least six. CASCADE only to a weak entity.
--     Write a one-line comment on each saying why.
\end{lstlisting}
\begin{lstlisting}[language=bash]
# --- 8. Create it ---------------------------------------------------
rm -f university.db
sqlite3 university.db ".read code/schema/v1_create.sql"
sqlite3 university.db ".tables"
sqlite3 university.db ".schema section"

# --- 9. Check it against the diagram, constraint by constraint ------
sqlite3 university.db "PRAGMA foreign_key_list(section);"
sqlite3 university.db "PRAGMA foreign_key_list(enrolment);"
sqlite3 university.db "PRAGMA table_info(course);"
# For EVERY line of Figure 6.1, find the line of SQL that implements
# it. Any line of the diagram with no SQL is a mapping you skipped.

# --- 10. Prove the M:N key does its job ------------------------------
sqlite3 university.db "PRAGMA foreign_keys=ON;
  INSERT INTO student VALUES ('2021-IT-001','Ayesha Khan','BSIT',2);
  INSERT INTO department VALUES ('CS','Computer Science');
  INSERT INTO course VALUES ('CS-103','Database Systems',4,'CS');
  INSERT INTO section VALUES ('CS-103','FALL2026',1,'R7',NULL);
  INSERT INTO enrolment VALUES
    ('2021-IT-001','CS-103','FALL2026',1,NULL);
  INSERT INTO enrolment VALUES
    ('2021-IT-001','CS-103','FALL2026',1,NULL);"
# The second one must fail. Record the message.
\end{lstlisting}

\smallskip
\textbf{What to notice.} Step 9 is the check that makes mapping a procedure
rather than a craft: every line of the diagram should have a line of SQL, and
every line of SQL should have a line of the diagram. Lines of SQL with no
diagram behind them are decisions you made without noticing, and they are
exactly the ones that turn out to be wrong in week twelve.
\end{lab}

\begin{checkpoint}
\begin{tightnum}
  \item A 1:1 relationship joins \ent{department} and \ent{instructor} (the
        head of department). Say which table gets the foreign key, justify the
        choice from the participation constraints, and say what the column
        would look like if you chose the other side.
  \item \ent{enrolment}'s primary key has four attributes. Explain what fact
        about the university each one contributes, and say precisely what
        becomes possible if you remove \texttt{semester} from the key.
  \item Map the recursive relationship \ent{course} \ent{requires}
        \ent{course}. Give the table, its primary key, its two foreign keys,
        and one \sqlkw{CHECK} constraint that the ER diagram does not state but
        that the university certainly means.
\end{tightnum}
\end{checkpoint}

\begin{chaptersummary}
\begin{tight}
  \item Mapping is \textbf{mechanical}. Seven rules, applied in order, and two
        people get the same schema --- which is what four chapters of diagram
        bought.
  \item \textbf{Entities first} (rules 1--3), then relationships (4--7): a
        foreign key can only point at a primary key that already exists.
  \item A \textbf{composite} attribute is stored as its parts; a
        \textbf{multi-valued} attribute becomes a table whose primary key is
        \emph{both} columns.
  \item \textbf{1:N} puts the foreign key on the \textbf{many} side, because
        one column holds one value. \textbf{1:1} puts it on the \emph{totally
        participating} side so it can be \sqlkw{NOT NULL}.
  \item \textbf{M:N} becomes a table carrying both keys plus the
        relationship's own attributes; the primary key is the combination, and
        it is what forbids enrolling twice.
  \item The M:N table is not plumbing. It holds the grades, it records the
        time, and it is the biggest table in the database.
  \item A \textbf{weak entity} becomes a table whose key is the owner's key
        plus the partial key, with the owner's key simultaneously a foreign
        key and \sqlkw{NOT NULL}.
  \item A \textbf{recursive} M:N uses rule 6 unchanged, taking the same key
        twice under the ER diagram's \textbf{role names}.
  \item A \textbf{ternary} relationship is one table with three foreign keys,
        never three binary relationships --- that loses the fact it existed to
        record.
  \item A multi-column foreign key needs its own \sqlkw{FOREIGN KEY} clause,
        with columns in the parent key's order.
  \item Check the schema back: every line of the diagram should have a line of
        SQL, and every line of SQL a line of the diagram.
\end{tight}
\end{chaptersummary}

\begin{reviewq}
  \item \textbf{(MCQ)} In a 1:N relationship the foreign key is placed:
        (a)~on the one side; (b)~on the many side; (c)~in a new table;
        (d)~on either side.
  \item \textbf{(MCQ)} A multi-valued attribute maps to: (a)~several columns;
        (b)~a comma-separated column; (c)~a new table keyed on owner and
        value; (d)~a \sqlkw{CHECK} constraint.
  \item \textbf{(Short)} Why can an M:N relationship not be mapped with a
        foreign key? Answer using a property of relations from
        \chref{relational}.
  \item \textbf{(Short)} A 1:1 relationship is partial on both sides. Where do
        you put the foreign key, and what is the consequence either way?
  \item \textbf{(Short)} What does the composite primary key of
        \ent{enrolment} forbid, and what would be permitted without it?
  \item \textbf{(Applied)} Map this diagram completely: a \ent{library} has
        many \ent{book}s; each book has many \ent{copy}s; \ent{member}s
        \ent{borrow} copies and a borrowing records a date out and a date
        back. Give every table, every key, every foreign key and a referential
        action for each, and name the rule that produced each table.
  \item \textbf{(Applied)} Take your own \texttt{v1\_create.sql} and go through
        Figure~\ref{fig:ercomplete} line by line. Produce a two-column table:
        the diagram constraint, and the SQL that implements it. Then list
        anything in either column with no partner, and for each one say
        whether the diagram or the schema is wrong.
\end{reviewq}
